% !TeX root = ../main.tex

% ==================================================
% A.9. 인수를 받는 데코레이터
% ==================================================

\subsection{인수를 받는 데코레이터와 데코레이터 팩토리}

Decorator Factory · Higher-Order Function · Nested Closure · Argument Binding · Configuration State

\subsubsection{이전 학습과의 연결}

\zsh{A.6}부터 \zsh{A.8}까지는 다음과 같은 데코레이터를 중심으로 학습했다.

\begin{dhkcode}
from functools import wraps

def decorator(func):

    @wraps(func)
    def wrapper(*args, **kwargs):
        return func(*args, **kwargs)

    return wrapper


@decorator
def add(a, b):
    return a + b
\end{dhkcode}

여기서 \zsh{@decorator}는 일반적인 단일 데코레이터 상황에서 다음 변환식과 대응했다.

\begin{dhkcode}
add = decorator(add)
\end{dhkcode}

이 구조에는 두 번의 함수 호출이 관여한다.

첫 번째는 함수 정의문이 실행될 때 발생하는 \zsh{decorator(add)} 호출이고, 두 번째는 이후 \zsh{add(10, 20)}을 실행할 때 발생하는 wrapper 호출이다.

\zsh{A.8}에서는 이 구조에 \zsh{functools.wraps}를 추가하여 원본 함수의 메타데이터와 \zsh{__wrapped__} 참조를 유지하는 방법까지 학습했다.

그런데 다음과 같은 기능을 구현하려면 어떻게 해야 할까?

\begin{dhkcode}
@repeat(3)
def greet():
    print("Hello")
\end{dhkcode}

이 코드는 \zsh{greet()}을 한 번 호출할 때 원본 함수의 본문을 세 번 실행하도록 만들려는 의도를 가지고 있다.

여기에는 두 종류의 인수가 존재한다.

\begin{dhksmry}
    - 3: 데코레이터의 동작을 설정하는 함수

    - greet()에 전달하는 인수: 실제 원본 함수를 실행할 때 사용하는 인수
\end{dhksmry}

이 둘은 서로 다른 함수 호출에 속한다.

이번 학습에서는 이 차이를 바탕으로 데코레이터를 생성하는 함수, 원본 함수를 전달받는 함수, 실제 호출을 중개하는 함수를 각각 구별할 것이다.

\subsubsection{이번 학습의 핵심 질문}

이번 학습에서 해결할 질문은 다음과 같다.

\begin{dhksmry}
1. @repeat(3)은 왜 @repeat와 다른가?

2. repeat(3)(greet)에는 어떤 함수 객체와 호출 결과가 관여하는가?

3. 세 단계로 중첩된 함수는 각각 어떤 네임스페이스를 가지는가?

4. 설정값 times와 원본 함수 객체 func는 어느 클로저 셀에 유지되는가?

5. repeat(2)와 repeat(5)로 만든 데코레이터는 상태를 공유하는가?

6. 데코레이터가 반환한 wrapper는 원본 함수의 인수와 반환값을 어떻게 
   처리하는가?

7. 데코레이터의 설정값을 변경하거나 검증하려면 어떤 조건을 고려해야 하는가?
\end{dhksmry}

가장 중요한 개념은 다음이다.

\begin{dhksmry}[centered]
인수를 받는 데코레이터는 일반적으로 데코레이터를 반환하는 함수,

즉 데코레이터 팩토리(Decorator Factory)를 통해 구현한다.
\end{dhksmry}

여기서 팩토리는 새로운 종류의 Python 객체가 아니다.

다른 함수를 생성하거나 반환하는 역할을 수행하는 일반적인 호출 가능한 객체를 의미한다.

\subsubsection{\zsh{@repeat(3)}의 기본 변환식}

먼저 두 문법을 비교해 보자.

\begin{dhknote}
\begin{dhksmry}[centered]
인수가 없는 데코레이터
\end{dhksmry}
\end{dhknote}

\begin{dhkcode}
@decorator
def greet():
    print("Hello")
\end{dhkcode}

대응 관계:

\begin{dhkcode}
greet = decorator(greet)
\end{dhkcode}

\begin{dhknote}
\begin{dhksmry}[centered]
인수를 받는 데코레이터
\end{dhksmry}
\end{dhknote}

\begin{dhkcode}
@repeat(3)
def greet():
    print("Hello")
\end{dhkcode}

대응 관계:

\begin{dhkcode}
greet = repeat(3)(greet)
\end{dhkcode}

두 번째 변환식의 오른쪽을 분석해 보자.

\begin{dhkcode}
repeat(3)(greet)
\end{dhkcode}

이 표현식은 하나의 함수에 인수 두 개를 전달하는 호출이 아니다.

다음과 같이 두 번의 호출이 연속으로 이루어지는 표현식이다.

\begin{dhkcode}
decorator = repeat(3)

replacement = decorator(greet)
\end{dhkcode}

첫 번째 호출인 \zsh{repeat(3)}의 결과가 호출 가능한 객체여야 두 번째 호출인
\zsh{decorator(greet)}이 가능하다.

일반적으로 이 결과가 원본 함수를 전달받는 데코레이터 함수 객체다.

이후 데코레이터가 wrapper를 반환한다면 다음과 같은 관계가 성립한다.

\begin{dhkbranch}[in=8]
node main :: repeat(3)
arrow main
node main :: decorator 함수 객체
line main
line main :: 원본 greet 함수 객체 전달
arrow main
node main :: decorator(greet)
arrow main
node main :: wrapper 함수 객체
arrow main
node main :: greet 이름에 바인딩
\end{dhkbranch}

그다음 사용자가 \zsh{greet()}을 호출하면 세 번째 함수 호출 단계가 시작된다.

\begin{dhkbranch}[in=8]
node main :: greet()
arrow main
node main :: wrapper()
branchline main -> a1 :: 원본 greet() 호출
stop a1
branchline main -> a2 :: 원본 greet() 호출
stop a2
handoff main -> a3 :: 원본 greet() 호출
\end{dhkbranch}

따라서 \zsh{@repeat(3)} 구조에서는 다음 세 호출을 구별해야 한다.

\begin{dhktable}
|단계|	호출 표현식	목적|	시점|
|:---|:---|:---|
|1|	repeat(3)|	설정값을 받아 데코레이터 생성|	정의문 실행 중|
|2|	decorator(//*greet)|	원본 함수를 받아 wrapper 생성|	정의문 실행 중|
|3|	greet()|	wrapper를 통해 원본 함수 실행|	이후 실제 호출 시|
\end{dhktable}

이 세 단계를 혼동하지 않는 것이 이번 학습의 출발점이다.

\subsubsection{\zsh{repeat(3)}을 직접 구현해 보자}

먼저 독립적으로 실행할 수 있는 전체 코드를 살펴본다.

\begin{dhkcode}
from functools import wraps

def repeat(times):
    print("factory:", times)

    def decorator(func):
        print("decorator:", func.__name__)

        @wraps(func)
        def wrapper(*args, **kwargs):
            print("wrapper start")

            result = None

            for i in range(times):
                print("iteration:", i + 1)
                result = func(*args, **kwargs)

            print("wrapper end")
            return result

        return wrapper

    return decorator


@repeat(3)
def greet(name):
    print("Hello", name)
    return name.upper()


print("definition complete")

print("result:", greet("Kim"))
\end{dhkcode}

예상 출력:

\begin{dhkzsh}
factory: 3
decorator: greet
definition complete
wrapper start
iteration: 1
Hello Kim
iteration: 2
Hello Kim
iteration: 3
Hello Kim
wrapper end
result: KIM
\end{dhkzsh}

이 출력은 세 함수의 실행 시점을 명확하게 보여준다.

\zsh{factory: 3}은 \zsh{repeat(3)}이 호출될 때 출력된다.

\zsh{decorator: greet}은 반환된 데코레이터가 원본 \zsh{greet}함수 객체를 전달받을 때 출력된다.

반면 \zsh{wrapper start}부터 \zsh{wrapper end}까지는 실제로 \zsh{greet("Kim")}을 호출한 뒤에 출력된다.

즉, 데코레이터를 생성하는 것, 원본 함수를 데코레이션하는 것, 데코레이션된 함수를 실행하는 것은 서로 다른 단계다.

이제 각 단계에서 어떤 객체가 생성되고 어떤 이름이 바인딩되는지 분석해 보자.

\subsubsection{첫 번째 단계 - \zsh{repeat(3)}과 설정값 바인딩}

다음 표현식이 평가된다.

\begin{dhkcode}
repeat(3)
\end{dhkcode}

이때 \zsh{repeat} 함수가 호출되고 매개변수 \zsh{times}가 정수 객체 3에 바인딩된다.

\begin{dhkcode}
def repeat(times):
    ...
\end{dhkcode}

함수 본문에서 다음 중첩 함수 정의문이 실행된다.

\begin{dhkcode}
def decorator(func):
    ...
\end{dhkcode}

여기서 새로운 \zsh{decorator} 함수 객체가 생성된다.

그다음 다음 문장을 만난다.

\begin{dhkcode}
return decorator
\end{dhkcode}

이 문장은 \zsh{decorator()}를 호출하는 것이 아니다.

생성된 \zsh{decorator} 함수 객체 자체를 반환하는 문장이다.

그렇다면 \zsh{repeat()} 호출이 종료된 뒤에도 \zsh{decorator}는 설정값 \zsh{times}를 사용할 수 있을까?

가능하다.

왜냐하면 \zsh{decorator} 내부에서 정의된 \zsh{wrapper}가 \zsh{times}를 참조하기 때문이다.

Python은 중첩된 함수들이 외부 함수의 바인딩에 접근할 수 있도록 클로저 환경을 구성한다.

여기서 \zsh{A.4}의 자유 변수와 셀 객체 개념이 다시 등장한다.

\begin{dhkdir}[in=6]
repeat(3) 호출
> times -> 3
> decorator 함수 객체 생성
>> times 바인딩에 접근 가능
\end{dhkdir}

\zsh{repeat()} 호출이 끝났다고 해서 \zsh{times}를 참조하는 클로저 환경까지 즉시 소멸하는 것은 아니다.

반환된 함수 객체가 해당 바인딩에 접근할 수 있는 상태를 유지하는 것이다.

\subsubsection{두 번째 단계 - \zsh{decorator(greet)}과 원본 함수 바인딩}

이제 \zsh{repeat(3)}이 반환한 함수 객체가 원본 \zsh{greet} 함수 객체를 인수로 전달받아 호출된다.

개념적으로 다음과 같다.

\begin{dhkcode}
configured_decorator = repeat(3)

replacement = configured_decorator(greet)
\end{dhkcode}

여기서 \zsh{configured_decorator}는 \zsh{repeat(3)}이 반환한 \zsh{decorator} 함수 객체를 참조한다.

그리고 \zsh{configured_decorator(greet)} 호출에서는 매개변수 \zsh{func}가 원본 \zsh{greet} 함수 객체에 바인딩된다.

\begin{dhkcode}
def decorator(func):

    @wraps(func)
    def wrapper(*args, **kwargs):
        ...
\end{dhkcode}

이 단계에서 새로운 \zsh{wrapper} 함수 객체가 생성된다.

Wrapper는 두 가지 외부 이름을 사용한다.

\begin{dhkzsh}
times
func
\end{dhkzsh}

\zsh{times}는 가장 바깥쪽 함수인 \zsh{repeat}의 매개변수이고, \zsh{func}는 바로 바깥쪽 함수인 \zsh{decorator}의 매개변수다.

따라서 wrapper는 서로 다른 렉시컬 스코프에 정의된 두 바인딩에 접근해야 한다.

\zsh{@wraps(func)}도 이 단계에서 적용된다.

\zsh{A.8}에서 학습한 것처럼 wrapper의 메타데이터가 설정되고 \zsh{__wrapped__} 속성이 원본 함수 객체를 참조하게 된다.

마지막으로 다음 문장이 실행된다.

\begin{dhkcode}
return wrapper
\end{dhkcode}

이제 데코레이터는 wrapper 함수 객체를 반환하고, 그 객체가 최종적으로 모듈 이름 \zsh{greet}에 바인딩된다.

\begin{dhkgraph}[uniform]
node n1 (1,1) :: 이름 greet
node n2 (1,2) :: wrapper 함수 객체
node n3 (2,1) :: Cell times // 3
node n4 (2,2) :: Cell func // 원본 함수 참조
node n5 (3,2) :: 원본 greet 함수 객체 // 실제 Hello 출력 코드

edge n1 -> n2
split n2 -> n3, n4
edge n4 -> n5

explain :: 데코레이션 이후

caption :: Wrapper는 설정값 times와 원본 함수 객체 func를 서로 다른 클로저 셀을 통해 참조한다. 그림은 핵심 바인딩 관계를 단순화한 것이다.
\end{dhkgraph}

여기서 중요한 점은 \zsh{times}와 \zsh{func}가 같은 함수의 매개변수가 아니라는 사실이다.

\zsh{times}는 \zsh{repeat()} 호출에서 바인딩되고, \zsh{func}는 \zsh{decorator()} 호출에서 바인딩된다.

두 이름이 만들어지는 시점과 네임스페이스가 서로 다르다.

\subsubsection{세 번째 단계 - Wrapper 호출과 반복 실행}

이제 다음 문장을 실행한다고 하자.

\begin{dhkcode}
greet("Kim")
\end{dhkcode}

현재 모듈 이름 \zsh{greet}은 wrapper 함수 객체를 참조한다.

따라서 wrapper가 호출되고 다음 바인딩이 형성된다.

\begin{dhkzsh}
args = ("Kim",)
kwargs = {}
\end{dhkzsh}

Wrapper는 자신의 클로저를 통해 \zsh{times}가 참조하는 정수 객체 3을 조회한다.

그다음 다음 반복문을 실행한다.

\begin{dhkcode}
for i in range(times):
    result = func(*args, **kwargs)
\end{dhkcode}

반복문은 세 번 실행된다.

각 반복에서 \zsh{func}가 참조하는 원본 \zsh{greet} 함수 객체가 호출된다.

원본 함수는 \zsh{"Hello Kim"}을 출력하고 문자열 \zsh{"KIM"}을 반환한다.

그 반환 객체가 wrapper의 지역 이름 \zsh{result}에 바인딩된다.

세 번의 호출이 완료되면 마지막 호출의 반환 객체가 \zsh{result}에 바인딩된 상태로 남아 있다.

따라서 다음 문장은

\begin{dhkcode}
return result
\end{dhkcode}

마지막 원본 함수 호출의 반환 객체를 최종 호출자에게 전달한다.

여기서 중요한 점이 있다.

\zsh{repeat(3)}이 원본 함수의 반환값을 세 배로 만드는 것은 아니다. 원본 함수를 세 번 호출하는 것이다.

원본 함수가 상태를 변경한다면 그 상태 변경도 세 번 발생할 수 있다.

\subsubsection{\zsh{repeat(3)}의 평가 시점과 함수 정의 시점}

\zsh{A.6}에서는 데코레이터가 함수 정의문 실행 중 적용된다고 설명했다.

이번에는 한 단계 더 정확하게 구별한다.

데코레이터 표현식 \zsh{repeat(3)}은 원본 함수 객체가 생성되기 전에 평가된다.

반면 \zsh{repeat(3)}이 반환한 데코레이터는 원본 함수 객체가 생성된 다음 적용된다.

다음 코드로 확인해 보자.

\begin{dhkcode}
def factory(times):
    print("factory:", times)

    def decorator(func):
        print("decorator:", func.__name__)
        return func

    return decorator

@factory(2)
def greet(value=print("default evaluated")):
    print("body")

print("definition complete")
greet()
\end{dhkcode}

예상 출력:

\begin{dhkzsh}
factory: 2
default evaluated
decorator: greet
definition complete
body
\end{dhkzsh}

왜 이런 순서일까?

먼저 \zsh{@factory(2)}라는 데코레이터 표현식이 평가된다.

따라서 \zsh{factory(2)}가 호출되고 데코레이터 함수 객체가 반환된다.

그다음 함수 정의에 필요한 기본 인수 표현식이 평가된다.

\begin{dhkcode}
value=print("default evaluated")
\end{dhkcode}

이 표현식은 \zsh{"default evaluated"}를 출력하고 \zsh{None}을 반환한다.

그 결과가 기본 인수 객체로 사용된다.

이후 원본 \zsh{greet} 함수 객체가 생성되고, 앞에서 준비한 데코레이터가 적용된다.

마지막으로 데코레이터의 반환 객체가 \zsh{greet}이라는 이름에 바인딩된다.

따라서 \zsh{A.9.3}의 변환식

\begin{dhkcode}
greet = repeat(3)(greet)
\end{dhkcode}

은 데코레이터 적용의 의미를 설명하는 데 유용하지만, 모든 부작용의 평가 순서까지 완전히 재현하는 문자 그대로의 소스 코드 치환은 아니다.

이 차이는 데코레이터 표현식이나 기본 인수 표현식에 부작용이 있을 때 중요해진다.

\subsubsection{세 단계의 클로저를 직접 검증하기}

이번에는 \zsh{A.4}와 \zsh{A.5}에서 사용한 \zsh{co_cellvars}, \zsh{co_freevars}, \zsh{__closure__}를 활용해 내부 구조를 확인해 보자.

\begin{dhkcode}
from functools import wraps

def repeat(times):

    def decorator(func):

        @wraps(func)
        def wrapper(*args, **kwargs):
            result = None

            for _ in range(times):
                result = func(*args, **kwargs)

            return result

        return wrapper

    return decorator


configured = repeat(2)


def square(x):
    return x * x


wrapped = configured(square)

print(repeat.__code__.co_cellvars)

print(configured.__code__.co_freevars)
print(configured.__code__.co_cellvars)

print(wrapped.__code__.co_freevars)

closure_map = dict(
    zip(wrapped.__code__.co_freevars, wrapped.__closure__)
)

print(closure_map["times"].cell_contents)
print(closure_map["func"].cell_contents is square)

print(wrapped.__wrapped__ is square)
\end{dhkcode}

예상 출력:

\begin{dhkzsh}
('times',)
('times',)
('func',)
('func', 'times')
2
True
True
\end{dhkzsh}

각 결과를 분석해 보자.

\begin{dhknote}
\begin{dhksmry}[centered]
1. repeat의 셀 변수
\end{dhksmry}
\end{dhknote}

\begin{dhkcode}
repeat.__code__.co_cellvars
\end{dhkcode}

결과는 다음과 같다.

\begin{dhkzsh}
('times',)
\end{dhkzsh}

\zsh{times}는 \zsh{repeat}의 지역 매개변수이지만 중첩 함수에서도 참조된다.

따라서 클로저 셀을 통해 유지할 필요가 있다.

\begin{dhknote}
\begin{dhksmry}[centered]
2. decorator의 자유 변수
\end{dhksmry}
\end{dhknote}

\begin{dhkcode}
configured.__code__.co_freevars
\end{dhkcode}

결과는 다음과 같다.

\begin{dhkzsh}
('times',)
\end{dhkzsh}

\zsh{decorator} 함수는 \zsh{repeat}의 렉시컬 환경에서 정의되었다.

그 내부에 정의되는 wrapper가 \zsh{times}를 사용하므로, 이 바인딩을 중첩된 함수까지 전달할 수 있는 클로저 구조가 형성된다.

\begin{dhknote}
\begin{dhksmry}[centered]
3. decorator의 셀 변수
\end{dhksmry}
\end{dhknote}

\begin{dhkcode}
configured.__code__.co_cellvars
\end{dhkcode}

결과는 다음과 같다.

\begin{dhkzsh}
('func',)
\end{dhkzsh}

\zsh{func}는 \zsh{decorator}의 지역 매개변수다.

그런데 wrapper가 이 이름을 참조하므로 셀 변수로 관리된다.

\begin{dhknote}
\begin{dhksmry}[centered]
4. wrapper의 자유 변수
\end{dhksmry}
\end{dhknote}

\begin{dhkcode}
wrapped.__code__.co_freevars
\end{dhkcode}

결과는 다음과 같다.

\begin{dhkzsh}
('func', 'times')
\end{dhkzsh}

Wrapper는 원본 함수 객체아ㅗ 설정값에 모두 접근해야 하므로 두 자유 변수를 사용한다.

이때 \zsh{__closure__}의 셀 순서를 임의로 가정하면 안 된다.

따라서 다음과 같이 이름과 셀을 대응시키는 것이 정확하다.

\begin{dhkcode}
closure_map = dict(
    zip(wrapped.__code__.co_freevars, wrapped.__closure__)
)
\end{dhkcode}

이후 다음 표현식으로 각 바인딩을 확인할 수 있다.

\begin{dhkcode}
closure_map["times"].cell_contents
closure_map["func"].cell_contents
\end{dhkcode}

여기서 핵심은 다음이다.

\begin{dhksmry}[centered]
데코레이터 팩토리는 중첩 함수가 세 단계라는 점이 특별한 것이 아니라,

서로 다른 단계에서 형성된 바인딩을 최종 wrapper가 함께 참조한다.
\end{dhksmry}

이 점이 중요하다.

\subsubsection{서로 다른 팩토리 호출은 상태를 공유할까?}

다음 코드를 살펴보자.

\begin{dhkcode}
from functools import wraps

def repeat(times):

    def decorator(func):

        @wraps(func)
        def wrapper(*args, **kwargs):
            result = None

            for _ in range(times):
                result = func(*args, **kwargs)

            return result

        return wrapper

    return decorator


@repeat(2)
def first():
    print("First")


@repeat(3)
def second():
    print("Second")


first()
second()
\end{dhkcode}

예상 출력:

\begin{dhkzsh}
First
First
Second
Second
Second
\end{dhkzsh}

여기서 \zsh{repeat()}이 두 번 호출된다.

첫 번째 호출에서는 \zsh{times}가 2를 참조하고, 두 번째 호출에서는 \zsh{times}가 3을 참조한다.

각 호출에서 독립적인 클로저 환경이 형성된다.

따라서 두 wrapper는 서로 다른 \zsh{times} 셀을 참조하게 된다.

이를 직접 검증할 수 있다.

\begin{dhkcode}
def cells_of(func):
    return dict(
        zip(func.__code__.co_freevars, func.__closure__)
    )


first_cells = cells_of(first)
second_cells = cells_of(second)

print(
    first_cells["times"] is second_cells["times"]
)
\end{dhkcode}

예상 출력:

\begin{dhkzsh}
False
\end{dhkzsh}

두 wrapper가 서로 다른 팩토리 호출에서 만들어진 셀을 참조하기 때문이다.

\begin{dhknote}
\begin{dhksmry}[centered]
반대로 하나의 데코레이터를 재사용한다면?
\end{dhksmry}
\end{dhknote}

다음 구조를 생각해 보자.

\begin{dhkcode}
configured = repeat(2)

@configured
def first():
    return 10


@configured
def second():
    return 20
\end{dhkcode}

여기서는 \zsh{repeat(2)}가 한 번만 호출된다.

반환된 동일한 \zsh{configured} 데코레이터 함수 객체를 두 함수에 적용한 것이다.

따라서 두 wrapper는 같은 \zsh{times} 셀을 공유할 수 있다.

하지만 각각 다른 원본 함수를 전달받았으므로 \zsh{func} 셀은 서로 다르다.

\begin{dhktable}
|비교 대상|	first와 second의 관계|
|:---|:---|
|times 셀|	동일한 셀을 공유|
|func 셀|	서로 다른 셀|
|wrapper 함수 객체|	서로 다른 객체|
|원본 함수 객체|	서로 다른 객체|
\end{dhktable}

이것은 팩토리 호출의 횟수와 데코레이터 적용 횟수가 서로 다르기 때문이다.

같은 팩토리 함수로 만들었다는 사실보다, 팩토리를 몇 번 호출했고, 어떤 반환 객체를 재사용했는지가 상태 공유 여부를 결정한다.

\subsubsection{설정값을 읽을 때 \zsh{nonlocal}이 필요할까?}

다음 코드를 보자.

\begin{dhkcode}
def repeat(times):

    def decorator(func):

        def wrapper(*args, **kwargs):
            for _ in range(times):
                func(*args, **kwargs)

        return wrapper

    return decorator
\end{dhkcode}

Wrapper는 \zsh{times}를 조회하지만 재바인딩하지 않는다.

따라서 \zsh{nonlocal times}가 필요하지 않다.

하지만 다음과 같이 변경하면 상황이 달라진다.

\begin{dhkcode}
def repeat(times):

    def decorator(func):

        def wrapper(*args, **kwargs):
            times += 1

            for _ in range(times):
                func(*args, **kwargs)

        return wrapper

    return decorator
\end{dhkcode}

이 코드는 wrapper를 호출할 때 \zsh{UnboundLocalError}가 발생할 수 있다.

왜냐하면 \zsh{times += 1}이 현재 함수 블록에서 일믕르 할당하는 연산이기 때문이다.

Python은 \zsh{times}를 wrapper의 지역 이름으로 분류하지만, 증가 연산을 위해 값을 읽으려는 시점에는 아직 지역 바인딩이 존재하지 않는다.

외부 \zsh{times}를 재바인딩하려면 다음과 같이 작성해야 한다.

\begin{dhkcode}
def wrapper(*args, **kwargs):
    nonlocal times

    times += 1
    ...
\end{dhkcode}

이 경우 \zsh{times}는 외부 \zsh{repeat} 호출의 바인딩을 변경하게 된다.

여기서 중요한 점은 \zsh{nonlocal}이 설정값을 읽기 위해 필요한 것이 아니라는 사실이다.

외부 바인딩을 재바인딩하려는 경우에 필요한 선언이다.

따라서 설정값을 읽기만 하는 일반적인 데코레이터 팩토리에서는 \zsh{nonlocal}을 사용하지 않아도 된다.

\subsubsection{설정값이 유효성 검사는 언제 수행해야 할까?}

지금까지 작성한 \zsh{repeat}에는 한 가지 문제가 있다.

다음 호출이 가능하다는 점이다.

\begin{dhkcode}
@repeat(0)
def greet():
    print("Hello")
\end{dhkcode}

기존 구현에서는 반복문이 한 번도 실행되지 않는다.

따라서 원본 함수도 호출되지 않고 \zsh{result}의 초기값인 \zsh{None}이 반환된다.

이것이 의도한 동작이라면 문제가 없지만, 반복 횟수가 반드시 양의 정수여야 한다는 계약이라면 유효성을 검사해야 한다.

\begin{dhkcode}
from functools import wraps


def repeat(times):

    if (
        not isinstance(times, int)
        or isinstance(times, bool)
        or times < 1
    ):
        raise ValueError("times는 양의 정수여야 합니다.")

    def decorator(func):

        @wraps(func)
        def wrapper(*args, **kwargs):
            result = None

            for _ in range(times):
                result = func(*args, **kwargs)

            return result

        return wrapper

    return decorator
\end{dhkcode}

이 검사는 \zsh{repeat(times)}가 호출되는 시점에 실행된다.

따라서 다음과 같은 코드는 함수 정의문 실행 중 오류가 발생한다.

\begin{dhkcode}
@repeat(0)
def greet():
    print("Hello")
\end{dhkcode}

왜냐하면 \zsh{repeat(0)}이 데코레이터 표현식으로 평가되는 순간 \zsh{ValueError}가 발생하기 때문이다.

여기서 \zsh{bool}을 별도로 검사한 이유도 있다.

Python에서는 다음 표현식이 \zsh{True}다.

\begin{dhkcode}
isinstance(True, int)
\end{dhkcode}

\zsh{bool}이 \zsh{int}의 하위 클래스이기 때문이다.

따라서 단순히 \zsh{isinstance(times, int)}만 검사하면 \zsh{True}도 정수로 받아들일 수 있다.

이처럼 설정값 검증은 데코레이터의 실행 의미를 명확하게 만드는 역할을 한다.

\subsubsection{\zsh{@repeat}과 \zsh{@repeat(3)}을 혼동하면 어떻게 될까?}

다음과 같이 정의된 팩토리를 생각해 보자.

\begin{dhkcode}
def repeat(times):

    def decorator(func):
        ...

    return decorator
\end{dhkcode}

이 함수는 원본 함수 객체가 아니라 설정값 \zsh{times}를 먼저 전달받도록 설계되었다.

따라서 정상적인 사용법은 다음과 같다.

\begin{dhkcode}
@repeat(3)
def greet():
    ...
\end{dhkcode}

반면 다음과 같이 작성하면 어떻게 될까?

\begin{dhkcode}
@repeat
def greet():
    ...
\end{dhkcode}

이때 Python은 개념적으로 다음 호출을 수행한다.

\begin{dhkcode}
greet = repeat(greet)
\end{dhkcode}

따라서 \zsh{times}에는 정수 3이 아니라 원본 \zsh{greet} 함수 객체가 전달된다.

이후 동작은 팩토리 구현에 따라 달라진다.

앞에서 작성한 유효성 검사가 있다면 함수 객체는 정수가 아니므로 정의 시점에
\zsh{ValueError}가 발생한다.

유효성 검사가 없다면 함수 객체가 \zsh{times}에 저장된 상태로 데코레이터 함수가 반환될 수 있고, 이후 호출에서 예상하지 못한 오류가 발생할 수 있다.

반대로 다음 코드도 주의해야 한다.

\begin{dhkcode}
@repeat()
def greet():
    ...
\end{dhkcode}

\zsh{repeat}이 필수 매개변수 \zsh{times}를 요구한다면 인수를 전달하지 않았으므로 \zsh{TypeError}가 발생한다.

정리하면 다음과 같다.

\begin{dhktable}
|문법|	실제 적용 방식|
|:---|:---|
|@repeat|	repeat(원본 함수 객체)|
|@repeat()|	repeat()의 반환값을 원본 함수에 적용|
|@repeat(3)|	repeat(3)의 반환값을 원본 함수에 적용|
|@repeat(times=3)|	repeat(times=3)의 반환값을 원본 함수에 적용|
\end{dhktable}

하나의 데코레이터를 \zsh{@repeat}과 \zsh{@repeat(3)} 양쪽 형태로 사용하도록 구현할 수도 있다.

하지만 이 경우에는 전달받은 객체가 설정값인지 원본 함수 객체인지 구별하는 추가 설계가 필요하다.

이번 학습에서는 두 형태를 명확하게 분리해서 이해하는 것이 우선이다.

\subsubsection{\zsh{return decorator}, \zsh{return wrapper}, \zsh{return wrapper()}의 차이}

데코레이터 팩토리를 처음 구현할 때 가장 자주 혼동하는 부분이다.

다음 구조를 다시 보자.

\begin{dhkcode}
def repeat(times):

    def decorator(func):

        def wrapper(*args, **kwargs):
            result = None

            for _ in range(times):
                result = func(*args, **kwargs)

            return result

        return wrapper

    return decorator
\end{dhkcode}

여기에는 서로 다른 세 반환 위치가 존재한다.

\begin{dhktable}
|위치|	반환 표현식|	반환 객체|
|:---|:---|:---|
|repeat|	return decorator|	데코레이터 함수 객체|
|decorator|	return wrapper|	Wrapper 함수 객체|
|wrapper|	return result|	원본 함수 호출의 결과 객체|
\end{dhktable}

이 중 하나라도 잘못 작성하면 호출 구조가 달라진다.

예를 들어 다음과 같이 작성했다고 하자.

\begin{dhkcode}
def decorator(func):

    def wrapper(*args, **kwargs):
        return func(*args, **kwargs)

    return wrapper()
\end{dhkcode}

여기서는 wrapper 함수 객체를 반환하는 것이 아니라, wrapper를 즉시 호출하고 그 결과를 반환한다.

따라서 원본 함수가 인수를 요구한다면 데코레이션 시점에 오류가 발생할 수 있다.

반대로 인수가 필요 없는 원본 함수라면 데코레이션 시점에 원본 함수가 실행되고, 그 반환 객체가 원래 함수 이름에 바인딩될 수 있다.

이것은 일반적인 wrapper 데코레이터에서 의도한 구조와 다르다.

따라서 다음 구분을 기억해야 한다.

\begin{dhkcode}
return decorator  # 데코레이터 함수 객체 반환
return wrapper    # wrapper 함수 객체 반환
return wrapper()  # wrapper를 호출한 결과 반환
\end{dhkcode}

\zsh{A.1}에서 학습한 함수 객체와 함수 호출 결과의 차이가 그대로 적용되는 것이다.

\subsubsection{실전 응용 - 재시도 횟수를 설정하는 데코레이터}

이번에는 단순 반복이 아니라 실제 프로그램에서 활용할 수 있는 재시도(Retry) 데코레이터를 구현해 보자.

함수 실행 중 일시적인 오류가 발생하면 지정된 횟수만큼 다시 호출하는 구조다.

\begin{dhkcode}
from functools import wraps

class TemporaryError(Exception):
    pass

def retry(max_attempts):

    if (
        not isinstance(max_attempts, int)
        or isinstance(max_attempts, bool)
        or max_attempts < 1
    ):
        raise ValueError("max_attempts는 양의 정수여야 합니다.")

    def decorator(func):

        @wraps(func)
        def wrapper(*args, **kwargs):

            for attempt in range(1, max_attempts + 1):

                try:
                    return func(*args, **kwargs)

                except TemporaryError:

                    if attempt == max_attempts:
                        raise

                    print(
                        f"재시도: {attempt}/{max_attempts}"
                    )

        return wrapper

    return decorator


state = {"calls": 0}


@retry(3)
def unstable_operation():

    state["calls"] += 1

    if state["calls"] < 3:
        raise TemporaryError("일시적인 실패")

    return "Success"


print(unstable_operation())
\end{dhkcode}

예상 출력:

\begin{dhkzsh}
재시도: 1/3
재시도: 2/3
Success
\end{dhkzsh}

이 코드를 순차적으로 분석해 보자.

\begin{dhknote}
\begin{dhksmry}[centered]
1. 데코레이터 팩토리 호출
\end{dhksmry}
\end{dhknote}

\begin{dhkcode}
retry(3)
\end{dhkcode}

\zsh{max_attempts}가 정수 객체 3에 바인딩된다.

유효성 검사가 통과하면 \zsh{decorator} 함수 객체가 반환된다.

\begin{dhknote}
\begin{dhksmry}[centered]
2. 원본 함수 데코레이션
\end{dhksmry}
\end{dhknote}

반환된 데코레이터가 원본 \zsh{unstable_operation} 함수 객체를 전달받아 호출된다.

\zsh{func}가 원본 함수 객체에 바인딩되고 wrapper 함수 객체가 생성된다.

Wrapper는 \zsh{max_attempts}와 \zsh{func}를 클로저를 통해 참조할 수 있다.

\begin{dhknote}
\begin{dhksmry}[centered]
3. 첫 번째 실제 호출
\end{dhksmry}
\end{dhknote}

\begin{dhkcode}
unstable_operation()
\end{dhkcode}

현재 이름은 wrapper를 참조하므로 wrapper가 실행된다.

반복문의 첫 번째 시도에서 원본 함수가 호출된다.

\zsh{state["calls"]}가 1이 되고 \zsh{TemporaryError}가 발생한다.

Wrapper는 해당 예외를 처리하고 다음 시도로 진행한다.

\begin{dhknote}
\begin{dhksmry}[centered]
4. 두 번째 실제 호출
\end{dhksmry}
\end{dhknote}

같은 wrapper 호출 안에서 원본 함수가 다시 호출된다.

\zsh{state["calls"]}가 2가 되고 \zsh{TemporaryError}가 발생한다.

아직 최대 시도 횟수에 도달하지 않았으므로 다음 반복으로 진행한다.

\begin{dhknote}
\begin{dhksmry}[centered]
5. 세 번째 실제 호출
\end{dhksmry}
\end{dhknote}

원본 함수가 세 번째로 호출된다.

이번에는 \zsh{state["calls"]}가 3이므로 예외가 발생하지 않고 \zsh{"Success"}를 반환한다.

따라서 다음 문장이 실행된다.

\begin{dhkcode}
return func(*args, **kwargs)
\end{dhkcode}

원본 함수의 반환 객체가 wrapper의 반환 객체가 되고, wrapper 실행이 종료된다.

\begin{dhknote}
\begin{dhksmry}[centered]
이 예제에서 중요한 설계 조건
\end{dhksmry}
\end{dhknote}

재시도 데코레이터는 원본 함수를 여러 번 호출할 수 있다.

따라서 원본 함수가 외부 시스템의 상태를 변경하는 경우에는 주의해야 한다.

예를 들어 결제 요청, 데이터베이스 갱신, 파일 생성처럼 부작용이 있는 연산은 실패 후 재호출할 때 중복 처리가 발생할 수 있다.

이러한 경우에는 멱등성(Idempotency), 트랜잭션, 중복 요청 식별 등의 조건을 함께 고려해야 한다.

또한 모든 예외를 무조건 재시도하는 것이 적절하지는 않다.

위 예제에서는 의도적으로 \zsh{TemporaryError}만 처리하도록 설계했다.

즉, 데코레이터의 설정값은 단순한 반복 횟수가 아니라 호출 정책을 정의하는 구성 요소가 될 수 있다.

\subsubsection{데코레이터 팩토리와 클래스 인스턴스의 상태 비교}

이전 학습에서 클로저와 클래스 인스턴스의 상태 유지 방식을 비교했다.

이번에는 설정값을 가진 데코레이터 팩토리에도 같은 관점을 적용해 본다.

\begin{dhkcode}
@repeat(3)
def greet():
    ...
\end{dhkcode}

여기서 반복 횟수 3은 클로저 셀을 통해 유지된다.

반면 클래스 기반으로 비슷한 긴으을 구현한다면 다음과 같은 구조를 생각할 수 있다.

\begin{dhkcode}
class Repeat:

    def __init__(self, times):
        self.times = times

    def __call__(self, func):

        def wrapper(*args, **kwargs):
            result = None

            for _ in range(self.times):
                result = func(*args, **kwargs)

            return result

        return wrapper
\end{dhkcode}

이 클래스의 인스턴스는 호출 가능한 객체다.

따라서 다음과 같이 사용할 수 있다.

\begin{dhkcode}
@Repeat(3)
def greet():
    print("Hello")
\end{dhkcode}

여기서 \zsh{@Repeat(3)}은 \zsh{Repeat} 인스턴스를 생성한다.

그리고 생성된 인스턴스가 데코레이터로 호출되면서 원본 함수 객체를 전달받는다.

함수 기반 팩토리와 클래스 기반 팩토리를 비교하면 다음과 같다.

\begin{dhktable}
|비교 기준|	함수 기반 팩토리|	클래스 기반 팩토리|
|:---|:---|:---|
|설정값 저장|	클로저 셀|	인스턴스 속성|
|설정값 접근|	자유 변수 times|	self.times|
|데코레이터 역할|	반환된 함수 객체|	__call__을 가진 인스턴스|
|원본 함수 보존|	wrapper의 클로저|	wrapper의 클로저 등|
|설정값 변경|	외부 바인딩 재바인딩 또는 가변 객체 변경|	인스턴스 속성 재바인딩|
\end{dhktable}

여기서 두 방식의 목적은 유사하지만 내부 상태 저장 구조는 다르다.

함수 기반 팩토리는 클로저 환경을 이용하고, 클래스 기반 팩토리는 인스턴스 네임스페이스를 이용한다.

또한 클래스 기반 구현에서 \zsh{self.times}를 변경하면 이미 생성된 \zsh{wrapper}의 동작에도 영향을 줄 수 있다.

왜냐하면 wrapper가 실행될 때ㅔ마다 동일한 인스턴스의 \zsh{self.times} 속성을 조회하기 때문이다.

이것은 \zsh{A.5}에서 학습한 클로저의 상태 변경과 클래스의 속성 변경이 서로 다른 메커니즘으로 유사한 동작을 구현할 수 있다는 사례다.

\subsubsection{Generator와 Coroutine에 적용할 때의 경계 조건}

\zsh{A.7}에서 Generator Function에 일반적인 wrapper를 적용하면 Generator Object 생성과 실제 본문 실행 시점이 달라질 수 있다는 사실을 학습했다.

이번 학습의 \zsh{repeat}에도 같은 문제가 존재한다.

\begin{dhkcode}
@repeat(3)
def numbers():
    yield 1
    yield 2
\end{dhkcode}

이 경우 wrapper가 원본 Generato Function을 세 번 호출하면 Generator Object가 세 개 생성될 수 있다.

하지만 각 Generator Object의 본문이 자동으로 실행되는 것은 아니다.

따라서 \zsh{repeat(3)}이 1, 2를 세 번 생성하는 반복 시퀀스를 자동으로 만들어 준다고 생각하면 안 된다.

마찬가지로 \zsh{async def}로 정의된 코루틴 함수에 동기 wrapper를 적용하면, 반복 호출에서 코루틴 객체만 여러 개 생성할 수 있다.

따라서 데코레이터의 설계는 원본 함수가 일반 함수인지, Generator Function인지, Coroutine Function인지에 따라 달라질 수 있다.

이 문제의 핵심은 함수 호출이 반환하는 객체와 함수 본문의 실제 실행 시점을 구별하는 것이다.

\subsubsection{자주 발생하는 오해와 판단 기준}

\begin{dhktable}
|오해	|정확한 설명|
|:---|:---|
|@repeat(3)은 원본 함수에 3을 전달한다.	|3은 데코레이터 팩토리의 설정값이다.|
|repeat(3)이 실행되면 원본 함수도 세 번 실행된다.|	일반적인 구현에서는 데코레이터가 생성될 뿐이다.|
|팩토리, 데코레이터, wrapper는 같은 함수 객체다.|	일반적으로 서로 다른 함수 객체다.|
|return decorator와 return decorator()는 같다.|	전자는 함수 객체 반환, 후자는 호출 결과 반환이다.|
|Wrapper가 times를 읽으려면 nonlocal이 필요하다.|	외부 바인딩을 읽는 데는 필요하지 않다.|
|같은 팩토리 함수를 사용하면 모든 설정값을 공유한다.	|각 팩토리 호출의 클로저 환경은 일반적으로 독립적이다.|
|@wraps(func)가 설정값을 저장한다.	|설정값은 일반적으로 클로저에 유지되고, wraps는 메타데이터를 설정한다.|
|반복 호출의 반환값은 자동으로 모두 합쳐진다.	|반환값 처리 방식은 wrapper 구현에 따라 결정된다.|
|재시도 데코레이터는 어떤 함수에도 안전하게 적용할 수 있다.	|부작용, 예외 종류, 멱등성 등을 고려해야 한다.|
\end{dhktable}

\subsubsection{핵심 정리}

이번 학습의 핵심 구조는 다음과 같다.

\begin{dhkcode}
from functools import wraps

def decorator_factory(config):

    def decorator(func):

        @wraps(func)
        def wrapper(*args, **kwargs):
            # config와 func를 활용
            return func(*args, **kwargs)

        return wrapper

    return decorator
\end{dhkcode}

이 코드는 세 단계의 함수 객체와 서로 다른 바인딩 환경으로 구성된다.

\begin{dhksmry}
데코레이터 팩토리의 전체 실행 구조

1. 팩토리 호출
   decorator_factory(config)가 실행되고 설정값이 외부 함수의 
   매개변수에 바인딩된다. 데코레이터 함수 객체가 생성되어 반환된다.

2. 데코레이터 적용
   반환된 데코레이터가 원본 함수 객체를 전달받아 호출된다. func가 
   원본 함수 객체에 바인딩되고 wrapper 함수 객체가 생성된다.

3. 이름 재바인딩
   데코레이터가 반환한 wrapper 함수 객체가 원래 함수 이름에 바인딩된다.

4. 실제 함수 호출
   이후 사용자가 데코레이션된 이름을 호출하면 wrapper가 실행된다. 
   Wrapper는 클로저를 통해 설정값과 원본 함수 객체에 접근한다.

5. 반환 또는 예외 전파
   Wrapper의 구현에 따라 원본 함수의 호출 횟수, 인수 전달,
   반환값 처리, 예외 전파가 결정된다.
\end{dhksmry}

이번 학습에서 가장 중요한 판단 기준은 다음이다.

\begin{dhksmry}[centered]
@factory(config)를 분석할 때는 함수가 세 단계로 중첩되어 있다는

표면적인 구조보다, 각 함수가 언제 호출되고 어떤 객체를 반환하며

어떤 바인딩을 클로저로 유지하는지를 추적해야 한다.
\end{dhksmry}

이 기준을 적용하면 반복, 재시도, 캐싱 정책, 로깅 수준 설정 등 다양한 데코레이터 팩토리를 동일한 객체 모델로 이해할 수 있다.
